% Corrigé intégré automatiquement pour la banque Exercices.
% Source : A_Integrer/DDS_Reserve/999_TEC_Clever/999_TEC_Clever.tex
\begin{corrigebox}[Corrigé]
\CorrigeQuestion{1}
On paramètre les mouvements relatifs, puis on applique la composition des vitesses et, si nécessaire, la dérivation vectorielle dans le repère adapté. Pour un train épicycloïdal, on écrit la relation de Willis avec les nombres de dents avant d'imposer les éléments bloqués ou entraînés.
\CorrigeQuestion{2}
On définit les variables, le pas de calcul et les conditions initiales, puis on traduit l'équation par une relation de récurrence. L'algorithme est vérifié sur les premiers pas et la stabilité numérique est contrôlée en comparant le pas aux constantes de temps du système.
\CorrigeQuestion{3}
On choisit le système isolé et le référentiel galiléen, puis on calcule le torseur cinétique et le torseur dynamique au point le plus commode. Le PFD s'écrit $\{\mathcal T_{ext}\}=\{\mathcal D\}$ ; ses équations scalaires donnent les efforts ou la loi de mouvement.
\CorrigeQuestion{4}
La résolution consiste à identifier les données et l'inconnue, choisir la loi du cours adaptée, écrire l'équation symbolique avant toute application numérique, puis vérifier l'homogénéité, le signe et l'ordre de grandeur du résultat. Les hypothèses utilisées doivent être explicitement contrôlées à la fin.
\end{corrigebox}
